% Part III-A, abstract (included by main_en.tex).
A policy rewarded by the tests it can see may learn to satisfy them instead of solving the problem. We measure this in
multi-turn code repair: Qwen3.5-4B-Base writes a function, sees the first failing example test, and tries again, up
to three times; held-out tests decide whether the final program is correct. Offline, \NUM{cs_hack_all} of the base
policy's attempts that pass the visible tests fail a held-out test, and both \kev{} and an autoregressive judge of the
same size rank such attempts below correct ones (AUROC \NUM{csauc_all_kevg} and \NUM{csauc_all_arg}) and separate
hard-coded lookups of the tests from reference solutions almost perfectly. We then train with GRPO for 100 steps, two
seeds per reward, on the visible tests alone, on the visible tests times \kev{}'s judgment that the code is genuine, and
on all tests as an oracle. We saw no sign that the visible-test reward was exploited: the share of problems solved rose on KodCode
validation (\NUM{cgszmean_kod_allpass}\% to \NUM{cgall_vis_kod}\%), MBPP+ and HumanEval+, the share of visible-test
passes that are wrong fell by about half on two of the three sets, and programs that copy values from the visible
tests stayed as rare as in the base policy. The judge factor changed nothing beyond the variation between seeds, and the
oracle reward, which never pays a wrong program, left slightly fewer wrong passes and solved no more problems. This version replaces earlier versions whose results
were not produced by experiments. Code, logs and every number in this paper are in the
repository.
